\hypertarget{appendix-241c0e42ac2d}{}\label{referencestart:appendix-241c0e42ac2d}
\section{Quantized matvec and dequantization}\label{appendix-241c0e42ac2d-quantized-matvec-and-dequantization}

\hypertarget{appendix-241c0e42ac2d-quantized-matvec-and-dequantization}{}

Quantization trades representation precision for storage and bandwidth. The runtime question is not merely ``how many bits per weight?'' It is what mathematical value each stored code represents, where its scale or zero-point metadata lives, and how the dot product is accumulated without silently changing the model.

For a simple symmetric per-row signed 8-bit scheme,

\[
s = \frac{\max_i |w_i|}{127},
\qquad
q_i = \operatorname{round}\!\left(\frac{w_i}{s}\right),
\]

with clamping to the representable range. The reconstructed value is

\[
\hat w_i = s q_i.
\]

A reference matrix-vector product can therefore compute

\[
y = \sum_i (s q_i)x_i.
\]

That scalar implementation is not necessarily fast, but it creates an oracle for SIMD or GPU fused dequantization-and-dot kernels.

Before clamping, the scalar reconstruction error for ordinary round-to-nearest is bounded by roughly half a quantization step. The output error is workload-dependent because the dot product combines weight errors with input magnitudes. An accuracy report therefore needs representative activations or a model-level metric, not only weight reconstruction error.

\subsection{Zero rows and exceptional values}\label{appendix-241c0e42ac2d-zero-rows-and-exceptional-values}

\hypertarget{appendix-241c0e42ac2d-zero-rows-and-exceptional-values}{}

When all weights in a block are zero, dividing by the maximum absolute value would produce an invalid scale. The format and runtime need an explicit zero-block representation. Non-finite input weights should be rejected or handled by a documented conversion policy instead of entering integer casts accidentally.

Groupwise 4-bit formats add another trade-off. Smaller groups let the scale adapt more locally but spend more bytes on metadata and may increase metadata bandwidth. Larger groups reduce metadata overhead but make each scale describe a wider distribution.

\subsection{Reference versus optimized execution}\label{appendix-241c0e42ac2d-reference-versus-optimized-execution}

\hypertarget{appendix-241c0e42ac2d-reference-versus-optimized-execution}{}

tools/labs/phase29\_ml\_runtime\_lab.cpp implements a checked per-row Q8 reference and compares its dot product against the original floating-point row under a stated error budget. The point is not to advertise Q8 as the final format; it is to establish the proof pattern:

stored codes plus metadata → scalar interpretation → reference result → optimized kernel → measured error and performance.

A Metal or SIMD kernel should be judged against that reference on adversarial widths, tails, zero blocks, extreme scales and random vectors. SIMD lane packing is an implementation detail; the dequantized mathematical contract remains the same.

Related: \href{../07-Vision-XR-and-Playable-Worlds-P02/main.pdf\#nameddest=appendix-fdc7879e8bf2}{Quantization and compressed weights (Vision, XR, and Playable Worlds, Part 2, Reference notes: Quantization and compressed weights)}, \hyperref[appendix-fef72dd246f2]{GGUF binary model loading}, \hyperref[appendix-cebadf3106d1]{Tensor runtime layout and views}.
